\documentclass[header={Lecture 11 Exam Study Notes}]{study-note}

\begin{document}

\StudyTitleBlock{11}{Lecture 11: Exam Study Notes}{Spectrum in Hilbert spaces, selfadjoint operators, and Hilbert--Schmidt theorem}

\vspace{0.8em}

\begin{focusbox}
This lecture specializes spectral theory to Hilbert spaces, especially compact and selfadjoint operators. For the exam, this is where the spectral theorem becomes concrete.

\medskip

The inner-product convention is the same as in Lecture 8: $(\cdot,\cdot)_H$ is linear in the first argument and conjugate-linear in the second.
\end{focusbox}

\tableofcontents

\section{Compact Operators On Hilbert Spaces}

\begin{thmbox}{Theorem 14.4}
\thmCompactSpectrum
\end{thmbox}

\begin{prooflevelbox}{STATEMENT}
Know the statement: nonzero spectral values of a compact operator are eigenvalues with finite-dimensional eigenspaces, and $0$ is the only possible accumulation point. Not the full proof.
\end{prooflevelbox}

\begin{focusbox}
For compact operators, $0$ is special: it may belong to the spectrum even when it is not an eigenvalue. The theorem only says that every \emph{nonzero} spectral value is an eigenvalue.
\end{focusbox}

\section{Selfadjoint Operators}

\begin{defbox}{Symmetric and selfadjoint}
For a bounded operator $T\in\mathcal L(H)$:
\begin{itemize}
\item $T$ is \textbf{symmetric} if
\[
(Tx,y)_H=(x,Ty)_H \qquad \forall x,y\in H;
\]
\item $T$ is \textbf{selfadjoint} if $T=T^*$.
\end{itemize}
For bounded everywhere-defined operators on a Hilbert space, these are equivalent.
\end{defbox}

\begin{defbox}{Selfadjoint operator}
An operator $T\in\mathcal L(H)$ is \textbf{selfadjoint} if
\[
T=T^*.
\]
\end{defbox}

\begin{thmbox}{Orthogonality of eigenspaces}
If $T$ is selfadjoint and
\[
Tx=\lambda x,
\qquad
Ty=\mu y,
\qquad
\lambda\neq \mu,
\]
then $(x,y)_H=0$.
\end{thmbox}

\begin{prooflevelbox}{FULL}
Reproduce the calculation with the convention linear in the first variable:
\[
\lambda(x,y)_H=(Tx,y)_H=(x,Ty)_H=\overline{\mu}(x,y)_H.
\]
Since eigenvalues of a selfadjoint operator are real by Theorem~\textbf{14.6}, this becomes $(\lambda-\mu)(x,y)_H=0$.
\end{prooflevelbox}

\paragraph{Reason.}
Since $T=T^*$,
\[
\lambda(x,y)_H=(Tx,y)_H=(x,Ty)_H=\overline{\mu}(x,y)_H.
\]
Eigenvalues of a selfadjoint operator are real, so $\overline{\mu}=\mu$. Hence $(\lambda-\mu)(x,y)_H=0$, and if $\lambda\neq\mu$ then $(x,y)_H=0$.

\begin{thmbox}{Theorem 14.5: Hellinger--Toeplitz theorem}
\thmHellingerToeplitz
\end{thmbox}

\begin{prooflevelbox}{SKETCH}
Know that symmetry implies the graph is closed, then apply the closed graph theorem.
\end{prooflevelbox}

\paragraph{Connection with Lecture 6.}
The proof is an application of the closed graph theorem: symmetry is used to show the graph is closed.

\begin{thmbox}{Theorem 14.6: Spectrum of a selfadjoint bounded operator}
\thmSelfadjointSpectrumReal
\end{thmbox}

\begin{prooflevelbox}{SKETCH}
Know the key steps: show $\sigma(T)\subset\R$ by proving $\lambda I-T$ is invertible when $\operatorname{Im}\lambda\neq 0$, and then locate the spectrum inside $[m,M]$ via the Rayleigh quotient.
\end{prooflevelbox}

\begin{exbox}{Precise interval}
If
\[
m:=\inf_{\norm{x}=1}(Tx,x)_H,
\qquad
M:=\sup_{\norm{x}=1}(Tx,x)_H,
\]
then
\[
\sigma(T)\subset [m,M].
\]
This is the infinite-dimensional analogue of the fact that eigenvalues of a real symmetric matrix are real and lie between the extremal Rayleigh quotients.
\end{exbox}

\section{Spectral Theorem For Compact Selfadjoint Operators}

\begin{thmbox}{Theorem 14.7: Spectral theorem for compact selfadjoint operators}
\thmHilbertSchmidt
\end{thmbox}

\begin{prooflevelbox}{STATEMENT}
Know the statement precisely and the representation $Tx=\sum_n\lambda_n(x,e_n)e_n$. Not the full proof, but understand the logical flow: extract eigenvalues via variational principles, diagonalize.
\end{prooflevelbox}

\paragraph{Meaning.}
In the separable case, compact selfadjoint operators admit a full diagonalization by a countable orthonormal basis. This is the infinite-dimensional analogue of diagonalizing a symmetric matrix.

\paragraph{Important details.}
The eigenvalues are real, the nonzero eigenspaces are finite-dimensional, eigenvectors for different eigenvalues are orthogonal, and the only possible accumulation point of the eigenvalues is $0$.

\begin{focusbox}
This theorem is the spectral theorem for compact selfadjoint operators in the form most useful for the exam.
\end{focusbox}

\section{What To Memorize Precisely}

\begin{focusbox}
Know precisely:
\begin{enumerate}
\item the statements of Theorems \textbf{14.4}, \textbf{14.5}, \textbf{14.6}, \textbf{14.7};
\item the difference between symmetric and selfadjoint, and why for bounded everywhere-defined operators they coincide;
\item orthogonality of eigenspaces for distinct eigenvalues of a selfadjoint operator;
\item that selfadjoint bounded operators have real spectrum;
\item that compact selfadjoint operators admit an orthonormal eigenbasis;
\item the representation
\[
Tx=\sum_n \lambda_n(x,e_n)e_n.
\]
\end{enumerate}
\end{focusbox}

\end{document}
